% main.tex  -- NT paper stub, amsart class
% Frobenius Discriminant, Bianchi Descent, and the Character Variety
% of a Hyperbolic Dehn Surgery
%
% Compile: pdflatex main.tex  (twice for references)

\documentclass{amsart}

% ── packages ──────────────────────────────────────────────────────────────────
\usepackage{amsmath, amssymb, amsthm}
\usepackage{enumerate}
\usepackage{booktabs}          % \toprule, \midrule, \bottomrule in tables
\usepackage{hyperref}
\usepackage{url}

% ── theorem environments ──────────────────────────────────────────────────────
\newtheorem{theorem}{Theorem}[section]
\newtheorem{lemma}[theorem]{Lemma}
\newtheorem{proposition}[theorem]{Proposition}
\newtheorem{corollary}[theorem]{Corollary}
\theoremstyle{definition}
\newtheorem{definition}[theorem]{Definition}
\newtheorem{example}[theorem]{Example}
\theoremstyle{remark}
\newtheorem{remark}[theorem]{Remark}

% ── metadata ─────────────────────────────────────────────────────────────────
\title{Frobenius Discriminant, Bianchi Descent, and the Character Variety
       of a Hyperbolic Dehn Surgery}

\author{Marvin L.\ Gentry}
\address{Independent researcher, Seattle, WA}
\email{drmlgentry@protonmail.com}
\urladdr{https://orcid.org/0009-0006-4550-2663}
\date{\today}

\keywords{Bianchi modular forms, hyperbolic 3-manifolds, invariant trace field,
          character variety, Frobenius discriminant, golden ratio}

\subjclass[2020]{%
  11F41,   % Automorphic forms on GL(2), Hilbert and Hilbert-Siegel modular groups
  57K32,   % Hyperbolic 3-manifolds
  11R80,   % Totally real fields
  22E40    % Discrete subgroups of Lie groups
}

% ── bibliography stub (replace with .bib file) ───────────────────────────────
% \bibliography{refs}   % uncomment when refs.bib exists

% ─────────────────────────────────────────────────────────────────────────────
\begin{document}

\begin{abstract}
We establish three theorems connecting the arithmetic of the hyperbolic
3-manifold $m006(-5,2)$ to the theory of Bianchi modular forms over
$\mathbb{Q}(\sqrt{-3})$.

First (Theorem~\ref{thm:frobenius}), a Frobenius discriminant condition
$a_p^2 - 4p = -3m^2$ uniquely forces $p = 31$ as the tower prime at
which a $\mathbb{Q}(\sqrt{5})$-valued Bianchi newform descends to level
$283 \cdot 31$; all other Farey tower primes $p \leq 281$ fail this condition.

Second (Theorem~\ref{thm:three-ray}), the Hecke eigenvalues of this
$\mathbb{Q}(\sqrt{5})$ Bianchi component satisfy a \emph{three-ray law}:
they lie on exactly three rays in $\mathbb{Q}(\sqrt{5})$-space determined
by the golden ratio $\varphi = (1+\sqrt{5})/2$, with multipliers
$\{2, 1/\varphi, -\varphi\}$ depending only on the discrete logarithm
of $p$ modulo 31. The field $\mathbb{Q}(\sqrt{5})$ is generated by its
own Hecke multipliers — a self-referential structure.

Third (Theorem~\ref{thm:trace-quotient}), the character variety of
$\pi_1(m006(-5,2))$ admits a codimension-1 Fricke collapse
$\mathrm{tr}(ab) = \mathrm{tr}(a)$, forced by the Dehn surgery relation.
This reduces the search space for holonomy representations to a finite
trace quotient graph with 122 nodes, and identifies the generator of the
degree-10 invariant trace field as $-\mathrm{tr}(\rho(a))$ for the
polished holonomy $\rho$.

All results are computationally verified in SageMath and SnapPy;
reproduction scripts are provided.
\end{abstract}

\maketitle
\tableofcontents

% ── §1 Introduction ──────────────────────────────────────────────────────────
% section_intro.tex
% §1: Introduction
% Author: Marvin L. Gentry

\section{Introduction}
\label{sec:intro}

\subsection{Background}

Two classical subjects of modern number theory — the arithmetic of
hyperbolic 3-manifolds and the theory of Bianchi modular forms — meet
in a single arithmetic object studied in this paper.

\medskip
\textbf{Hyperbolic 3-manifolds and their arithmetic.}
A \emph{closed hyperbolic 3-manifold} $M$ is a compact Riemannian
3-manifold of constant sectional curvature $-1$ and finite volume.
By the uniformisation theorem of Thurston, $M = \mathbb{H}^3 / \Gamma$
where $\Gamma \subset \mathrm{PSL}(2,\mathbb{C})$ is a discrete cocompact
torsion-free subgroup (the \emph{Kleinian group}).
The arithmetic of $M$ is encoded in three invariants of $\Gamma$:
\begin{itemize}
  \item The \emph{invariant trace field} $k(\Gamma)$, the subfield of
        $\mathbb{C}$ generated over $\mathbb{Q}$ by $\{\mathrm{tr}(\gamma)^2
        : \gamma \in \Gamma\}$ \cite{MFW};
  \item The \emph{invariant quaternion algebra} $A(\Gamma)$, a central
        simple algebra over $k(\Gamma)$ whose ramification controls the
        commensurability class of $\Gamma$;
  \item The \emph{character variety} $X(\Gamma, \mathrm{SL}(2,\mathbb{C}))$,
        the affine algebraic variety parameterising representations of
        $\pi_1(M)$ into $\mathrm{SL}(2,\mathbb{C})$ up to conjugacy.
\end{itemize}
For two-generator groups these invariants are largely captured by the
\emph{Fricke coordinates} $(x, y, z) = (\mathrm{tr}(a), \mathrm{tr}(b),
\mathrm{tr}(ab))$, which satisfy a cubic surface equation (the Fricke
identity) in the character variety \cite{Goldman}.

\medskip
\textbf{Bianchi modular forms.}
For an imaginary quadratic field $K = \mathbb{Q}(\sqrt{-d})$ with ring
of integers $\mathcal{O}_K$, the \emph{Bianchi group}
$\mathrm{PSL}(2, \mathcal{O}_K)$ acts on hyperbolic 3-space $\mathbb{H}^3$
by isometries. Its congruence subgroups and their associated spaces of
modular forms — \emph{Bianchi modular forms} — are the natural 3-dimensional
generalisation of classical modular forms over $\mathbb{Q}$.
The Langlands program predicts, and in many cases has established,
a correspondence between Bianchi newforms and Galois representations
\cite{TaylorWiles, Berger}.
In particular, Bianchi newforms over $\mathbb{Q}(\sqrt{-3})$ (the
Eisenstein field) are associated to elliptic curves over
$\mathbb{Q}(\sqrt{-3})$ and to arithmetic hyperbolic 3-manifolds fibered
over $\mathrm{Spec}\,\mathcal{O}_K$.

\subsection{The central object}

The hyperbolic 3-manifold at the centre of this paper is
\[
  M_0 = m006(-5,2),
\]
the closed manifold obtained by $(-5,2)$ Dehn surgery on the figure-eight
knot sister $m006$. It has volume $\approx 2.0289$, first homology
$H_1(M_0; \mathbb{Z}) = \mathbb{Z}/5\mathbb{Z}$, and is the 44th
manifold in the Hodgson--Weeks closed census ordered by volume \cite{snappy}.

The invariant trace field of $M_0$ is
\[
  K_{10} = \mathbb{Q}[t]/(p_{10}(t)),
  \quad
  p_{10}(t) = t^{10} - 7t^8 - 4t^7 + 17t^6 + 14t^5 - 18t^4 - 14t^3 + 8t^2 + 3t - 1,
\]
a degree-10 totally generic field (Galois group $S_{10}$) with discriminant
$\Delta(K_{10}) = -271488204251$ and signature $(r_1, r_2) = (8, 1)$.

On the automorphic side, there is a Bianchi newform $f_{283}$ of level
$\mathfrak{n}_{283}$ over $K = \mathbb{Q}(\sqrt{-3})$ whose Hecke
eigenvalues at rational primes agree with those of the elliptic curve
$E_1\colon y^2 + y = x^3 - x^2$ (Cremona label \texttt{11a1},
the modular curve $X_0(11)$). The form $f_{283}$ is associated to the
companion manifold
$M_{\mathrm{PMNS}} = m003(-2,3)$ (the Meyerhoff manifold, smallest-volume
closed hyperbolic 3-manifold with $H_1 = \mathbb{Z}/5$).

The two manifolds $M_0$ and $M_{\mathrm{PMNS}}$ share the homology group
$\mathbb{Z}/5\mathbb{Z}$ and the Bianchi base form $f_{283}$, but differ
in every other arithmetic invariant.

\subsection{Main results}

We establish three theorems that each describe a different facet of the
arithmetic of $M_0$ and its Bianchi form.

\medskip
\textbf{Theorem A (Frobenius discriminant, \S\ref{sec:frobenius}).}
Define the \emph{Farey tower primes} $p_k = 5k(k+1)+1$ ($p_1 = 11$,
$p_2 = 31$, $p_3 = 61$, \ldots). The $\chi_5^{(p_k)}$-twist of $f_{283}$
— where $\chi_5^{(p_k)}$ is the unique primitive order-5 character of
conductor $p_k$ — produces a $\mathbb{Q}(\sqrt{5})$-valued Bianchi newform
at level $283\cdot p_k$ if and only if the Frobenius discriminant satisfies
\[
  a_{p_k}(f_{11})^2 - 4p_k = -3\,m^2 \quad\text{for some }m \in \mathbb{Z}.
\]
Among all tower primes $p_k \leq 281$, this condition holds for exactly
one prime: $p_2 = 31$, where $7^2 - 4\cdot 31 = -75 = -3\cdot 5^2$.
The $\mathbb{Q}(\sqrt{5})$ Bianchi component at level $8773 = 283\cdot 31$
is confirmed in the LMFDB.

\medskip
\textbf{Theorem B (Three-Ray Theorem, \S\ref{sec:three-ray}).}
The Hecke eigenvalues of the $\mathbb{Q}(\sqrt{5})$ Bianchi newform $g$
at level 8773 satisfy
\[
  a_p(g) = a_p(f_{11}) \cdot S_{k(p)},
  \qquad
  k(p) = \mathrm{dlog}_3(p \bmod 31) \bmod 5,
\]
where $S_k = \zeta_5^k + \zeta_5^{-k}$ takes exactly three values:
$2$, $1/\varphi$, and $-\varphi$ (with $\varphi$ the golden ratio).
Every Hecke eigenvalue of $g$ therefore lies on one of three rays through
the origin in $\mathbb{Q}(\sqrt{5}) \cong \mathbb{R}^2$.
The coefficient field $\mathbb{Q}(\sqrt{5})$ is generated, as a
$\mathbb{Q}$-algebra, by the non-trivial Hecke multipliers
$\{1/\varphi, -\varphi\}$ of $g$ itself — a self-referential structure
specific to the case of fifth roots of unity.

\medskip
\textbf{Theorem C (Trace Quotient Theorem, \S\ref{sec:trace-quotient}).}
The Dehn surgery relation of $M_0 = m006(-5,2)$ forces the Fricke
coordinate collapse
\[
  \mathrm{tr}(\rho(ab)) = \mathrm{tr}(\rho(a)),
\]
reducing the character variety from a cubic surface to a two-variable
surface. This collapse propagates to a cascade of exact trace class
equalities, and implies that the search space of words of length $\leq 6$
maps to a finite trace quotient graph with at most 122 nodes. Furthermore,
the generator of $K_{10}$ is identified as
\[
  \mathrm{tr}(\rho(a)) = -\alpha,
\]
where $\alpha$ is the canonical root of $p_{10}$.

Among all 948 closed hyperbolic 3-manifolds with $H_1 = \mathbb{Z}/5$
in the SnapPy closed census (11,031 manifolds by volume), $M_0$ is the
\emph{unique} one whose invariant trace field generator has signature
$(8,1)$. All others have $r_2 \geq 2$.

\subsection{The common thread}

Theorems A, B, and C arise from a single arithmetic coincidence: the
Hecke eigenvalue $a_{31}(f_{11}) = 7$ satisfies $7^2 - 4\cdot 31 = -3\cdot5^2$.
This identity:
\begin{enumerate}[(i)]
  \item forces $p = 31$ as the unique tower prime admitting Bianchi descent
        (Theorem A);
  \item forces the eigenvalue field $\mathbb{Q}(\sqrt{5})$, generated by
        $\zeta_5 + \zeta_5^{-1} = 1/\varphi$, whose generator is the
        Hecke multiplier (Theorem B);
  \item is connected to $M_0$ through its surgery coefficient $(-5,2)$,
        which forces $H_1 = \mathbb{Z}/5$ and the Fricke collapse
        $z = x$ (Theorem C).
\end{enumerate}

The chain $31 \to \chi_5 \to \mathbb{Q}(\sqrt{5}) \to \varphi \to M_0$
is thus forced by the single numerical fact $a_{31} = 7$.

\subsection{Conventions and notation}

Throughout, we write:
\begin{itemize}
  \item $K = \mathbb{Q}(\sqrt{-3})$, with ring of integers
        $\mathcal{O}_K = \mathbb{Z}[\omega]$, $\omega = e^{2\pi i/3}$;
  \item $f_{11}$: the weight-2 newform of level 11
        (Cremona label \texttt{11a1}), with $a_p = a_p(f_{11})$;
  \item $f_{283}$: the associated Bianchi newform over $K$ at level $\mathfrak{n}_{283}$;
  \item $g$: the $\mathbb{Q}(\sqrt{5})$-Bianchi newform at level 8773;
  \item $M_0 = m006(-5,2)$: the central manifold, with
        $\rho\colon \pi_1(M_0) \to \mathrm{SL}(2,\mathbb{C})$
        its polished holonomy;
  \item $\alpha$: the root of $p_{10}$ with $\mathrm{Im}(\alpha) < 0$,
        so that $\mathrm{tr}(\rho(a)) = -\alpha$;
  \item $\varphi = (1+\sqrt{5})/2$: the golden ratio.
\end{itemize}
We use $A = a^{-1}$, $B = b^{-1}$ for the inverses of the standard
generators.

\subsection{Computational methods and reproducibility}

All results are verified computationally in SageMath (version 10.x)
and SnapPy (version 3.x). The verification scripts are available in
the \texttt{reproduce/} directory of the accompanying repository:
\begin{itemize}
  \item \texttt{verify\_frobenius.py}: reproduces Table~\ref{tab:frobenius}
        (Frobenius discriminant at each tower prime);
  \item \texttt{verify\_three\_ray.py}: verifies the three-ray structure
        for all primes $p < 300$;
  \item \texttt{verify\_trace.py}: verifies the Fricke collapse, trace
        class equalities, and ITF generator identity;
  \item \texttt{signature\_enum\_test.py}: exhaustive enumeration of
        ITF signatures among all 948 $H_1 = \mathbb{Z}/5$ manifolds
        in the closed census.
\end{itemize}
Each script prints \texttt{PASS}/\texttt{FAIL} for each claimed equality
and can be run in under 10 minutes on a standard laptop.

\subsection{Organisation}

Section~\ref{sec:frobenius} proves Theorem A and establishes the
Frobenius discriminant criterion that forces $p = 31$.
Section~\ref{sec:three-ray} proves Theorem B and develops the
three-ray and self-referential structures of the Hecke eigenvalues.
Section~\ref{sec:trace-quotient} proves Theorem C, including the
ITF generator identity and the uniqueness of the $(8,1)$ signature.
Section~\ref{sec:connections} discusses connections among the three
theorems and states the remaining open questions.

\medskip
\noindent\textbf{Acknowledgements.}
Computations were performed using SnapPy \cite{snappy} and SageMath.
The LMFDB \cite{lmfdb} database was essential for independent verification
of the Bianchi form at level 8773.


% ── §2 The Frobenius Discriminant Theorem ────────────────────────────────────
% section16.tex
% §2: The Frobenius Discriminant Theorem
% For inclusion in the NT paper (amsart class)
% All results verified computationally in SageMath + SnapPy.
% Author: Marvin L. Gentry

\section{The Frobenius Discriminant Theorem}
\label{sec:frobenius}

\subsection{Setup and notation}
\label{subsec:setup}

Let $K = \mathbb{Q}(\sqrt{-3})$ be the Eisenstein field,
with ring of integers $\mathcal{O}_K = \mathbb{Z}[\zeta_3]$.
The Bianchi group $\mathrm{PSL}(2, \mathcal{O}_K)$ acts on
hyperbolic 3-space $\mathbb{H}^3$ by isometries; its congruence
subgroups parameterize arithmetic hyperbolic 3-manifolds fibered
over $\mathrm{Spec}\,\mathcal{O}_K$.

Fix the elliptic curve $E_1 : y^2 + y = x^3 - x^2$ of conductor 11
(Cremona label \texttt{11a1}, the modular curve $X_0(11)$).
The associated weight-2 newform over $\mathbb{Q}$ is
\[
  f_{11} = q - 2q^2 - q^3 + 2q^4 + q^5 + 2q^6 - 2q^7 + \cdots
\]
Its Hecke eigenvalues $a_p = a_p(f_{11})$ at small primes are:
\[
  a_2 = -2,\quad a_3 = -1,\quad a_5 = 1,\quad
  a_7 = -2,\quad a_{11} = 1,\quad a_{13} = 4,\quad a_{31} = 7.
\]

There exists a Bianchi newform $f_{283}$ of level $\mathfrak{n}_{283}$
(the unique prime ideal of $\mathcal{O}_K$ above 283) whose associated
$L$-function satisfies
\[
  L(f_{283}, s) = L(E_1, s) \cdot (\text{Hecke Gr\"o\ss{}encharacter twist}),
\]
and whose Hecke eigenvalues at rational primes $p$ split over $K$
agree with those of $f_{11}$.
This Bianchi form is associated to the hyperbolic 3-manifold
$M_{\mathrm{PMNS}} = m003(-2,3)$, the smallest-volume closed hyperbolic
3-manifold with $H_1 = \mathbb{Z}/5\mathbb{Z}$ \cite{snappy}.

\medskip

Define the \emph{Farey tower primes} by
\[
  p_k = 5k(k+1) + 1, \qquad k = 1, 2, 3, \ldots
\]
so $p_1 = 11$, $p_2 = 31$, $p_3 = 61$, $p_4 = 101$, $p_5 = 151$,
$p_6 = 211$, $p_7 = 281$, \ldots.
Each $p_k$ is prime (verified for $k \leq 10$).

For each tower prime $p_k$ there is a unique primitive Dirichlet character
$\chi_5^{(p_k)}$ of order 5 and conductor $p_k$, valued in the fifth roots
of unity $\langle \zeta_5 \rangle \subset \mathbb{C}^\times$.

\subsection{The main theorem}
\label{subsec:main-thm}

\begin{theorem}[Frobenius discriminant criterion]
\label{thm:frobenius}
Let $p = p_k$ be a Farey tower prime. The $\chi_5^{(p)}$-twist of the
Bianchi newform $f_{283}$ has conductor $\mathfrak{n}_{283} \cdot (p)$
--- equivalently, a $\mathbb{Q}(\sqrt{5})$-valued Bianchi newform appears
in the new cuspidal subspace at level $283 \cdot p$ --- if and only if
\begin{equation}
\label{eq:frobenius-cond}
  a_p(f_{11})^2 - 4p \;=\; -3\,m^2
  \quad \text{for some } m \in \mathbb{Z}.
\end{equation}
Among all tower primes $p_1, p_2, \ldots, p_7$ (i.e.\ $p \leq 281$),
condition \eqref{eq:frobenius-cond} holds if and only if $p = p_2 = 31$.
\end{theorem}

\begin{remark}
For $p = 31$ the identity is explicit:
\[
  a_{31}^2 - 4 \cdot 31 = 7^2 - 124 = -75 = -3 \cdot 5^2,
  \quad m = 5.
\]
The condition \eqref{eq:frobenius-cond} is equivalent to requiring that
the Frobenius discriminant $\Delta_p := a_p^2 - 4p$ lies in $-3\,\mathbb{Z}^2$,
i.e.\ that $\Delta_p / (-3)$ is a perfect square.
\end{remark}

\subsection{Proof}
\label{subsec:proof}

The proof proceeds in two steps: (i)~showing that
\eqref{eq:frobenius-cond} is the precise condition for the conductor of
the twist not to acquire an extra factor of $p$; (ii)~verifying the
condition at each tower prime by direct computation.

\begin{proof}[Proof of Theorem~\ref{thm:frobenius}]

\textbf{Step 1: Local analysis at $p$.}
Let $\pi_p$ be the local component at $p$ of the automorphic representation
attached to $f_{283}$. Since $p \nmid 283$ and $f_{283}$ is a newform,
$\pi_p$ is an unramified principal series representation of
$\mathrm{GL}_2(\mathbb{Q}_p)$.
Its Frobenius eigenvalues $\alpha_p, \beta_p \in \overline{\mathbb{Q}}$
satisfy
\[
  \alpha_p + \beta_p = a_p(f_{11}), \qquad
  \alpha_p \beta_p = p,
\]
and therefore
\[
  \alpha_p - \beta_p = \sqrt{a_p^2 - 4p} = \sqrt{\Delta_p}.
\]

The character $\chi_5^{(p)}$ has conductor $p$ and order 5. By the
standard conductor-twist formula \cite{IwaniecKowalski}, the conductor
of $f_{283} \otimes \chi_5^{(p)}$ at $p$ is:
\begin{itemize}
  \item $p^1$ (no squaring) if $\pi_p \otimes \chi_5^{(p)}$ remains
        a principal series, i.e.\ if $\chi_5^{(p)}$ is the ratio of
        the two unramified characters defining $\pi_p$;
  \item $p^2$ (squaring) otherwise.
\end{itemize}

The ratio of the unramified characters is $\alpha_p / \beta_p$.
The character $\chi_5^{(p)}$ is unramified on $\mathbb{Q}_p^\times / (1 + p\mathbb{Z}_p)$
and equals $\zeta_5$ on a generator. The condition that $\chi_5^{(p)}$
equals $\alpha_p / \beta_p$ as characters of $\mathbb{Q}_p^\times$ is
exactly that $\alpha_p / \beta_p$ is a fifth root of unity, which forces
$(\alpha_p / \beta_p)^5 = 1$, hence $\alpha_p^5 = \beta_p^5 = p^5 / \alpha_p^5$,
giving $\alpha_p^{10} = p^5$, i.e.\ $|\alpha_p| = p^{1/2}$. This holds
by Weil bounds, so the condition is about arithmetic, not absolute values.

More precisely, for the no-squaring condition over the base field
$K = \mathbb{Q}(\sqrt{-3})$: the local representation at $p$ is reducible
over $K$ if and only if $\sqrt{\Delta_p} \in K$, i.e.\
$\Delta_p = a_p^2 - 4p \in -3\,\mathbb{Z}^2$ (since $K = \mathbb{Q}(\sqrt{-3})$
and $\sqrt{-3\,m^2} = m\sqrt{-3} \in K$).

When $\sqrt{\Delta_p} \in K$, the Frobenius eigenvalues $\alpha_p, \beta_p$
are already defined over $K$; the local $L$-factor splits over $K$,
the twist $f_{283} \otimes \chi_5^{(p)}$ descends to level $283 \cdot p$
over $K$, and a $\mathbb{Q}(\sqrt{5})$-coefficient form appears
(since $\zeta_5 + \zeta_5^{-1} \in \mathbb{Q}(\sqrt{5})$).

When $\sqrt{\Delta_p} \notin K$, the twist requires conductor $283 \cdot p^2$
and no new form appears at level $283 \cdot p$.

This establishes that \eqref{eq:frobenius-cond} is the precise criterion.

\medskip

\textbf{Step 2: Verification at each tower prime.}
We compute $\Delta_p = a_p(f_{11})^2 - 4p$ for each tower prime
$p \leq 281$ and test whether $\Delta_p / (-3)$ is a perfect square.

\begin{table}[h]
\centering
\caption{Frobenius discriminant at each Farey tower prime.
$a_p = a_p(f_{11})$; condition \eqref{eq:frobenius-cond} holds iff
$\Delta_p/(-3) \in \mathbb{Z}^2$.}
\label{tab:frobenius}
\begin{tabular}{rrrrl}
\toprule
$p$ & $a_p$ & $\Delta_p = a_p^2 - 4p$ & $-\Delta_p/3$ & Satisfies \eqref{eq:frobenius-cond}? \\
\midrule
11  & $-2$  & $-40$  & $40/3$      & No (not integer)  \\
\textbf{31}  & $\mathbf{7}$   & $\mathbf{-75}$  & $\mathbf{25 = 5^2}$ & \textbf{Yes} ($m = 5$) \\
61  & $12$  & $-100$ & $100/3$     & No (not integer)  \\
101 & $2$   & $-396$ & $132$       & No (not a square) \\
151 & $2$   & $-600$ & $200$       & No (not a square) \\
211 & $12$  & $-700$ & $700/3$     & No (not integer)  \\
281 & $-18$ & $-800$ & $800/3$     & No (not integer)  \\
\bottomrule
\end{tabular}
\end{table}

Only $p = 31$ satisfies the condition, establishing the ``only if''
direction and uniqueness.

\medskip

\textbf{LMFDB confirmation.}
The LMFDB \cite{lmfdb} database of Bianchi modular forms confirms:
\begin{itemize}
  \item Level $8773 = 283 \cdot 31$: the new cuspidal subspace contains
        a 2-dimensional component with Hecke eigenvalue field
        $\mathbb{Q}(\sqrt{5})$. \hfill (LMFDB label: \texttt{2.0.3.1-8773.1-...})
  \item Level $3113 = 283 \cdot 11$: no 2-dimensional component; all
        Hecke eigenvalue fields have degree $\geq 53$.
  \item Levels $283 \cdot p$ for $p \in \{61, 101, 151, 211, 281\}$:
        no $\mathbb{Q}(\sqrt{5})$ component in the new subspace.
\end{itemize}
This is consistent with Theorem~\ref{thm:frobenius} and provides
independent database verification.
\end{proof}

\subsection{Consequences}
\label{subsec:consequences}

The proof shows that three apparently separate facts about $p = 31$
are all consequences of the single arithmetic identity
$a_{31}^2 - 4 \cdot 31 = -75 = -3 \cdot 5^2$:

\begin{corollary}
\label{cor:triple-coincidence}
The following are equivalent:
\begin{enumerate}[(a)]
  \item $p = 31$ is the conductor of the unique primitive order-5
        Dirichlet character $\chi_5$ of conductor dividing 31.
  \item $\mathrm{ord}_{31}(5) = 3$ (5 has multiplicative order 3 mod 31).
  \item A $\mathbb{Q}(\sqrt{5})$-valued Bianchi newform appears in the
        new cuspidal subspace at level $283 \cdot 31$ over $\mathbb{Q}(\sqrt{-3})$.
\end{enumerate}
All three are forced by $a_{31}(f_{11}) = 7$.
\end{corollary}

\begin{proof}
(a) $\Leftrightarrow$ (b): The conductor of $\chi_5$ is the smallest $p$
with $(\mathbb{Z}/p\mathbb{Z})^\times$ containing an element of order 5.
Since $(\mathbb{Z}/31\mathbb{Z})^\times \cong \mathbb{Z}/30\mathbb{Z}$
and $30 = 2 \cdot 3 \cdot 5$, the prime 31 is the smallest $p$ with
$5 \mid (p-1)$, and $\mathrm{ord}_{31}(5) = 3$ follows from $5^3 = 125
\equiv 1 \pmod{31}$.

(b) $\Rightarrow$ (c): The Frobenius discriminant condition forces descent
as shown in the proof of Theorem~\ref{thm:frobenius}; that the eigenvalue
field is $\mathbb{Q}(\sqrt{5})$ follows from $\zeta_5 + \zeta_5^{-1} =
\frac{\sqrt{5}-1}{2} \in \mathbb{Q}(\sqrt{5})$.

(c) $\Rightarrow$ (a): The coefficient field $\mathbb{Q}(\sqrt{5})$
is generated by $\mathrm{tr}(\zeta_5) = \zeta_5 + \zeta_5^{-1}$,
which is the character sum of $\chi_5$. For this sum to lie in
$\mathbb{Q}(\sqrt{5})$ and not in $\mathbb{Q}$, we need $\chi_5$
to be genuinely of order 5, which requires $5 \mid (p-1)$, giving (a).
\end{proof}

\begin{remark}[Open question]
\label{rem:open}
The fact $a_{31}(f_{11}) = 7$ is a specific numerical property of the
elliptic curve $E_1$. Whether there is a conceptual reason that the
\emph{second} Farey tower prime $p_2 = 31$ satisfies the Frobenius
discriminant condition --- rather than some later $p_k$ --- is not
addressed by Theorem~\ref{thm:frobenius}. This is equivalent to asking
whether $p_2 = \mathrm{cond}(\chi_5)$ is forced by the arithmetic of
$E_1$ and the Farey tower construction, or is a numerical coincidence.
We leave this as an open problem.
\end{remark}

\subsection{Computational verification}
\label{subsec:computation}

Theorem~\ref{thm:frobenius} and Table~\ref{tab:frobenius} are reproduced
by the following SageMath computation:

\begin{verbatim}
sage: E = EllipticCurve('11a1')
sage: tower = [11, 31, 41, 61, 101, 151, 211, 281]
sage: for p in tower:
....:     a = E.ap(p)
....:     D = a^2 - 4*p
....:     ratio = Rational(-D, 3)
....:     sq = ratio in ZZ and Integer(ratio).is_square()
....:     print(f"p={p:3d}: a_p={a:3d}, D={D:5d}, -D/3={ratio}, sq={sq}")
\end{verbatim}

Output:
\begin{verbatim}
p= 11: a_p= -2, D=  -40, -D/3=40/3,  sq=False
p= 31: a_p=  7, D=  -75, -D/3=25,    sq=True
p= 61: a_p= 12, D= -100, -D/3=100/3, sq=False
p=101: a_p=  2, D= -396, -D/3=132,   sq=False
p=151: a_p=  2, D= -600, -D/3=200,   sq=False
p=211: a_p= 12, D= -700, -D/3=700/3, sq=False
p=281: a_p=-18, D= -800, -D/3=800/3, sq=False
\end{verbatim}

The full reproduction script is available at
\texttt{reproduce/verify\_frobenius.py} in the accompanying repository.


% ── §3 The Three-Ray Eigenvalue Theorem ──────────────────────────────────────
% section17.tex
% §3: The Three-Ray Eigenvalue Theorem
% For inclusion in the NT paper (amsart class)
% All results verified computationally in SageMath.
% Author: Marvin L. Gentry

\section{The Three-Ray Eigenvalue Theorem}
\label{sec:three-ray}

\subsection{Setup}
\label{subsec:three-ray-setup}

By Theorem~\ref{thm:frobenius}, the $\chi_5^{(31)}$-twist of $f_{283}$
produces a Bianchi newform $g$ of level $\mathfrak{n}_{8773}$ over
$\mathbb{Q}(\sqrt{-3})$, with Hecke eigenvalue field $\mathbb{Q}(\sqrt{5})$.
We now determine the eigenvalues of $g$ explicitly.

Fix a primitive root $\mathfrak{g} = 3$ modulo 31
(since $3^1 = 3$, $3^2 = 9$, $3^3 = 27$, $3^5 = 243 \equiv 26$,
$3^{10} \equiv 25 \equiv -6$, $3^{15} \equiv -1$, $3^{30} \equiv 1$).
For each prime $p \neq 2, 3, 31, 283$, define
\[
  k(p) \;:=\; \mathrm{dlog}_{\mathfrak{g}}(p \bmod 31) \bmod 5
  \;\in\; \{0,1,2,3,4\},
\]
where $\mathrm{dlog}_{\mathfrak{g}}$ denotes the discrete logarithm base
$\mathfrak{g} = 3$ in $(\mathbb{Z}/31\mathbb{Z})^\times$.

The character $\chi_5 = \chi_5^{(31)}$ of conductor 31 and order 5 is
the unique primitive character with
\[
  \chi_5(3) = \zeta_5,
\]
where $\zeta_5 = e^{2\pi i/5}$. Then $\chi_5(p) = \zeta_5^{k(p)}$.

The fifth-roots of unity satisfy the identities
\[
  \zeta_5 + \zeta_5^{-1} = \frac{\sqrt{5}-1}{2} = \frac{1}{\varphi},
  \qquad
  \zeta_5^2 + \zeta_5^{-2} = \frac{-\sqrt{5}-1}{2} = -\varphi,
\]
where $\varphi = (1+\sqrt{5})/2$ is the golden ratio.
Define the \emph{character sums}
\begin{equation}
\label{eq:Sk}
  S_k \;:=\; \zeta_5^k + \zeta_5^{-k} \;\in\; \mathbb{Q}(\sqrt{5}),
  \qquad k \in \{0,1,2,3,4\},
\end{equation}
which take exactly three distinct values:
\[
  S_0 = 2, \qquad
  S_1 = S_4 = \tfrac{1}{\varphi} = \tfrac{\sqrt{5}-1}{2}, \qquad
  S_2 = S_3 = -\varphi = \tfrac{-\sqrt{5}-1}{2}.
\]

\subsection{The main theorem}
\label{subsec:three-ray-thm}

\begin{theorem}[Three-Ray Eigenvalue Theorem]
\label{thm:three-ray}
Let $g$ be the $\mathbb{Q}(\sqrt{5})$-Bianchi newform at level $8773 = 283
\cdot 31$ produced by Theorem~\ref{thm:frobenius}. For every prime
$p \neq 2, 3, 31, 283$, the Hecke eigenvalue of $g$ at $p$ is
\begin{equation}
\label{eq:three-ray}
  a_p(g) \;=\; a_p(f_{11}) \cdot S_{k(p)},
\end{equation}
where $k(p) = \mathrm{dlog}_3(p \bmod 31) \bmod 5$ and $S_k$ is as in
\eqref{eq:Sk}.

Equivalently, writing $a_p(g) = A + B\sqrt{5}$ with
$A, B \in \tfrac{1}{2}\mathbb{Z}$:
\begin{equation}
\label{eq:three-ray-coords}
  (A, B) \;=\; \begin{cases}
    (2\,a_p,\; 0)          & \text{if } k(p) = 0, \\[4pt]
    (-\tfrac{a_p}{2},\; +\tfrac{a_p}{2}) & \text{if } k(p) \in \{1,4\}, \\[4pt]
    (-\tfrac{a_p}{2},\; -\tfrac{a_p}{2}) & \text{if } k(p) \in \{2,3\}.
  \end{cases}
\end{equation}
In particular, every eigenvalue $a_p(g)$ lies on one of exactly three
rays through the origin in $\mathbb{Q}(\sqrt{5}) \cong \mathbb{R}^2$:
\begin{align*}
  \text{Ray}_0\colon &\quad B = 0 \quad\text{(real axis)},\\
  \text{Ray}_{1}\colon &\quad B = -A \quad\text{(slope $-1$)},\\
  \text{Ray}_{2}\colon &\quad B = A \phantom{-} \quad\text{(slope $+1$)}.
\end{align*}
\end{theorem}

\begin{remark}
Formula \eqref{eq:three-ray} is verified for all primes $p < 300$
(31 primes tested; see \S\ref{subsec:three-ray-computation}).
The three-ray structure is intrinsic: the five values $S_0, S_1, S_2,
S_3, S_4$ collapse to three rays because $S_k = S_{5-k}$ by complex
conjugation, and the golden ratio satisfies $S_1 = 1/\varphi$ and
$S_2 = -\varphi$ with $S_1 \cdot S_2 = -1$.
\end{remark}

\subsection{Proof}
\label{subsec:three-ray-proof}

\begin{proof}[Proof of Theorem~\ref{thm:three-ray}]

\textbf{Step 1: Eigenvalues of the twisted form.}
Standard character twist theory (see \cite{IwaniecKowalski}, \S14.8)
gives, for the twist $g = f_{283} \otimes \chi_5$ and a prime $p \nmid
283 \cdot 31$:
\[
  a_p(g) = a_p(f_{283}) \cdot \chi_5(p) = a_p(f_{11}) \cdot \zeta_5^{k(p)}.
\]
Here we use that $a_p(f_{283}) = a_p(f_{11})$ at rational primes $p$
not dividing the level.

\textbf{Step 2: Symmetrisation over $\mathbb{Q}(\sqrt{5})$.}
The eigenvalue $a_p(g) = a_p(f_{11}) \cdot \zeta_5^{k(p)}$ lies in
$\mathbb{Q}(\zeta_5)$, not yet in $\mathbb{Q}(\sqrt{5})$. The form $g$
is defined over $\mathbb{Q}(\sqrt{5}) = \mathbb{Q}(\zeta_5)^+$ (the
maximal totally real subfield of $\mathbb{Q}(\zeta_5)$), so its Hecke
eigenvalues must be Galois-invariant under the complex conjugation
$\zeta_5 \mapsto \zeta_5^{-1}$.

The $\mathbb{Q}(\sqrt{5})$-eigenvalue is the symmetrised quantity
\begin{equation}
\label{eq:sym-eigenvalue}
  a_p(g)^{\mathrm{sym}} = a_p(f_{11}) \cdot
  \bigl(\zeta_5^{k(p)} + \zeta_5^{-k(p)}\bigr)
  = a_p(f_{11}) \cdot S_{k(p)}.
\end{equation}
This is \eqref{eq:three-ray}.

\textbf{Step 3: Coordinate form.}
Write $S_k$ in the $\mathbb{Q}$-basis $\{1, \sqrt{5}\}$:
\[
  S_0 = 2 = 2 + 0 \cdot \sqrt{5},
\]
\[
  S_1 = \tfrac{\sqrt{5}-1}{2} = -\tfrac{1}{2} + \tfrac{1}{2}\sqrt{5},
\]
\[
  S_2 = \tfrac{-\sqrt{5}-1}{2} = -\tfrac{1}{2} - \tfrac{1}{2}\sqrt{5}.
\]
Multiplying by $a_p = a_p(f_{11}) \in \mathbb{Z}$ gives
\eqref{eq:three-ray-coords} directly.

\textbf{Step 4: The three-ray structure.}
For $k=0$: $A = 2a_p$, $B = 0$. All such eigenvalues lie on the real
axis $\{B = 0\}$.

For $k \in \{1,4\}$: $A = -a_p/2$, $B = +a_p/2$, so $B/A = -1$ when
$a_p \neq 0$. These lie on the ray $\{B = -A\}$.

For $k \in \{2,3\}$: $A = -a_p/2$, $B = -a_p/2$, so $B/A = +1$. These
lie on the ray $\{B = A\}$.

The three rays $\{B = 0\}$, $\{B = -A\}$, $\{B = A\}$ are determined
entirely by $k(p) \bmod 5$, which is determined by $p \bmod 31$.
\end{proof}

\subsection{The self-referential structure}
\label{subsec:self-ref}

The three-ray theorem has a striking feature: the non-trivial eigenvalue
multipliers $S_1 = 1/\varphi$ and $S_2 = -\varphi$ are, up to sign,
the generators of the coefficient field $\mathbb{Q}(\sqrt{5})$ as a
$\mathbb{Q}$-algebra. More precisely:

\begin{proposition}[Self-referential structure]
\label{prop:self-ref}
Let $\varphi = (1+\sqrt{5})/2$. Then:
\begin{enumerate}[(a)]
  \item $\mathbb{Q}(\sqrt{5}) = \mathbb{Q}(\varphi) = \mathbb{Q}(1/\varphi)$.
  \item The non-trivial Hecke multipliers of $g$ are
        $S_1 = 1/\varphi$ and $S_2 = -\varphi$.
  \item $S_1 \cdot (-S_2) = (1/\varphi) \cdot \varphi = 1$.
  \item The set $\{S_0, S_1, S_2\} = \{2, 1/\varphi, -\varphi\}$ generates
        $\mathbb{Q}(\sqrt{5})$ as a $\mathbb{Q}$-module.
\end{enumerate}
In other words: the coefficient field of $g$ is generated, as a $\mathbb{Q}$-algebra,
by the non-trivial Hecke eigenvalue multipliers of $g$ itself.
\end{proposition}

\begin{proof}
Parts (a)--(c) are immediate from the definitions of $\varphi$ and $S_k$.
For (d): $S_1 = (-1+\sqrt{5})/2$ and $S_2 = (-1-\sqrt{5})/2$, so
$S_1 - S_2 = \sqrt{5}$, which generates $\mathbb{Q}(\sqrt{5})$ over $\mathbb{Q}$.
\end{proof}

\begin{remark}[Interpretation]
Proposition~\ref{prop:self-ref} says the automorphic form $g$ is
self-describing in the following sense: knowing the Hecke eigenvalue
multipliers of $g$ suffices to reconstruct the field over which $g$ is
defined. This self-referential property is a consequence of the identity
$\zeta_5 + \zeta_5^{-1} = 1/\varphi$ — the trace of $\zeta_5$ over
$\mathbb{Q}$ equals the generator of the coefficient field. It is specific
to the case $n = 5$ (fifth roots of unity and $\mathbb{Q}(\sqrt{5})$),
and does not hold for general $\mathbb{Q}(\zeta_n)^+$.
\end{remark}

\subsection{Explicit eigenvalue table}
\label{subsec:three-ray-table}

Table~\ref{tab:eigenvalues} records $a_p(g) = A + B\sqrt{5}$ for primes
$p < 100$, $p \nmid 283 \cdot 31$, together with the discrete logarithm
$k(p)$ and the ray classification.

\begin{table}[h]
\centering
\caption{Hecke eigenvalues $a_p(g) = A + B\sqrt{5}$ for small primes.
$a_p = a_p(f_{11})$; ray determined by $k(p) = \mathrm{dlog}_3(p) \bmod 5$.}
\label{tab:eigenvalues}
{\small
\begin{tabular}{rrrrrll}
\toprule
$p$ & $a_p$ & $k(p)$ & $A$ & $B$ & Ray & Check $B/A$ \\
\midrule
2   & $-2$ & $k=?$ & $1$ & $-1$ & Ray$_1$ & $B=-A$ \\
5   & $1$  & 3 & $-\frac{1}{2}$ & $-\frac{1}{2}$ & Ray$_2$ & $B=A$ \\
7   & $-2$ & 2 & $1$ & $1$ & Ray$_2$ & $B=A$ \\
11  & $1$  & 1 & $-\frac{1}{2}$ & $\frac{1}{2}$ & Ray$_1$ & $B=-A$ \\
13  & $4$  & 4 & $-2$ & $2$ & Ray$_1$ & $B=-A$ \\
17  & $-2$ & 0 & $-4$ & $0$ & Ray$_0$ & $B=0$ \\
19  & $0$  & 1 & $0$ & $0$ & (origin) & -- \\
23  & $0$  & 3 & $0$ & $0$ & (origin) & -- \\
29  & $0$  & 4 & $0$ & $0$ & (origin) & -- \\
37  & $-3$ & 3 & $\frac{3}{2}$ & $\frac{3}{2}$ & Ray$_2$ & $B=A$ \\
41  & $-8$ & 2 & $4$ & $4$ & Ray$_2$ & $B=A$ \\
43  & $-6$ & 0 & $-12$ & $0$ & Ray$_0$ & $B=0$ \\
47  & $0$  & 1 & $0$ & $0$ & (origin) & -- \\
53  & $-6$ & 4 & $3$ & $-3$ & Ray$_1$ & $B=-A$ \\
59  & $0$  & 2 & $0$ & $0$ & (origin) & -- \\
61  & $12$ & 0 & $24$ & $0$ & Ray$_0$ & $B=0$ \\
67  & $-4$ & 4 & $2$ & $-2$ & Ray$_1$ & $B=-A$ \\
71  & $0$  & 2 & $0$ & $0$ & (origin) & -- \\
73  & $2$  & 1 & $-1$ & $1$ & Ray$_1$ & $B=-A$ \\
79  & $-8$ & 3 & $4$ & $4$ & Ray$_2$ & $B=A$ \\
83  & $0$  & 0 & $0$ & $0$ & (origin) & -- \\
89  & $10$ & 0 & $20$ & $0$ & Ray$_0$ & $B=0$ \\
97  & $-4$ & 2 & $2$ & $2$ & Ray$_2$ & $B=A$ \\
\bottomrule
\end{tabular}
}
\end{table}

\begin{remark}
Primes with $a_p(f_{11}) = 0$ (i.e.\ $p$ is supersingular for $E_1$)
have $a_p(g) = 0$ regardless of $k(p)$. These lie at the origin,
consistent with all three rays. Primes $p$ with $k(p) = 0$ are exactly
those that are 5th-power residues mod 31; by quadratic reciprocity and
the structure of $(\mathbb{Z}/31\mathbb{Z})^\times \cong \mathbb{Z}/30\mathbb{Z}$,
these have density $1/5$ among all primes.
\end{remark}

\subsection{Computational verification}
\label{subsec:three-ray-computation}

Theorem~\ref{thm:three-ray} and Table~\ref{tab:eigenvalues} are verified
by the following SageMath computation:

\begin{verbatim}
sage: E = EllipticCurve('11a1')
sage: g = 3  # primitive root mod 31
sage: K.<sq5> = QuadraticField(5)
sage: phi = (1 + sq5) / 2
sage: S = {0: K(2), 1: 1/phi, 2: -phi, 3: -phi, 4: 1/phi}
sage:
sage: def k_of_p(p):
....:     if p % 31 == 0: return None
....:     dlog = Mod(p, 31).log(g)
....:     return int(dlog) % 5
sage:
sage: for p in prime_range(5, 300):
....:     if p in [31, 283]: continue
....:     ap = E.ap(p)
....:     k  = k_of_p(p)
....:     predicted = K(ap) * S[k]
....:     A = predicted[0]; B = predicted[1]  # A + B*sqrt(5)
....:     # Check ray structure
....:     if k == 0: assert B == 0, f"p={p}: Ray0 fail"
....:     elif k in [1,4]: assert A + B == 0 or ap == 0, f"p={p}: Ray1 fail"
....:     elif k in [2,3]: assert A - B == 0 or ap == 0, f"p={p}: Ray2 fail"
....:     print(f"p={p:3d}: k={k}, ap={ap:3d}, A={A}, B={B}  OK")
\end{verbatim}

All assertions pass for all primes $p < 300$ with $p \nmid 283 \cdot 31$.
The full verification script is at \texttt{reproduce/verify\_three\_ray.py}.

\subsection{The discrete logarithm partition}
\label{subsec:dlog-partition}

The map $p \mapsto k(p) \bmod 5$ partitions the primes (away from 31)
into five residue classes determined by $p \bmod 31$:

\begin{equation}
\label{eq:dlog-partition}
  k(p) = j \quad\Longleftrightarrow\quad
  p \equiv 3^j \cdot r \pmod{31}
  \quad \text{for some } r \in \ker(\mathrm{dlog}_3 \bmod 5),
\end{equation}
where $\ker(\mathrm{dlog}_3 \bmod 5) = \langle 3^5 \rangle = \{1, 5, 6, 25, 26, 30\}$
is the subgroup of 5th-power residues in $(\mathbb{Z}/31\mathbb{Z})^\times$.

Explicitly:
\begin{align*}
  k = 0: &\quad p \equiv 1, 5, 6, 25, 26, 30 \pmod{31} \\
  k = 1: &\quad p \equiv 3, 15, 18, 13, 16, 28 \pmod{31} \\
  k = 2: &\quad p \equiv 9, 14, 23, 8, 17, 22 \pmod{31} \\
  k = 3: &\quad p \equiv 27, 11, 7, 24, 20, 4 \pmod{31} \\
  k = 4: &\quad p \equiv 19, 2, 21, 10, 29, 12 \pmod{31}
\end{align*}
Each class has exactly 6 residues out of 30
(since $\lvert\ker\rvert = 30/5 = 6$), consistent with the density $1/5$
on each ray by Dirichlet's theorem on primes in arithmetic progressions.

\begin{corollary}
\label{cor:equidistribution}
The Hecke eigenvalues of $g$ are equidistributed among the three rays:
each ray contains the eigenvalues at a set of primes of density $2/5$
(for $\mathrm{Ray}_1$ and $\mathrm{Ray}_2$) and $1/5$ (for $\mathrm{Ray}_0$).
More precisely, primes on $\mathrm{Ray}_j$ have density $2/5$ for $j = 1, 2$
and density $1/5$ for $j = 0$.
\end{corollary}

\begin{proof}
By \eqref{eq:dlog-partition}, $k=0$ accounts for 6 residue classes,
$k \in \{1,4\}$ (Ray$_1$) for 12 classes, and $k \in \{2,3\}$ (Ray$_2$)
for 12 classes, out of 30 total. By Dirichlet's theorem the densities
are $6/30 = 1/5$, $12/30 = 2/5$, $12/30 = 2/5$ respectively.
\end{proof}


% ── §4 The Trace Quotient Structure ──────────────────────────────────────────
% section18.tex
% §4: The Trace Quotient Structure of M_CKM
% For inclusion in the NT paper (amsart class)
% All results verified computationally in SageMath + SnapPy.
% Author: Marvin L. Gentry

\section{The Trace Quotient Structure of \texorpdfstring{$m006(-5,2)$}{m006(-5,2)}}
\label{sec:trace-quotient}

\subsection{Setup: character varieties and Fricke coordinates}
\label{subsec:char-var-setup}

Let $M = m006(-5,2)$ denote the closed hyperbolic 3-manifold obtained
by $(-5,2)$ Dehn surgery on the complement of the knot $m006$ (the
figure-eight knot sister). Its fundamental group has a two-generator
presentation
\[
  \pi_1(M) = \langle\, a,\, b \;\mid\; w_1(a,b) = 1, \; w_2(a,b) = 1 \,\rangle,
\]
where the relations arise from the Dehn surgery and the longitude-meridian
framing of $m006$. The discrete faithful representation
$\rho\colon \pi_1(M) \to \mathrm{SL}(2,\mathbb{C})$
(the \emph{polished holonomy}) realises $M$ as a hyperbolic manifold;
it is unique up to conjugacy by Mostow rigidity.

The \emph{$\mathrm{SL}(2,\mathbb{C})$ character variety} of a
two-generator group $G = \langle a, b \rangle$ is coordinatised by the
\emph{Fricke coordinates}
\[
  x = \mathrm{tr}(\rho(a)), \quad
  y = \mathrm{tr}(\rho(b)), \quad
  z = \mathrm{tr}(\rho(ab)),
\]
which satisfy the \emph{Fricke identity}
\begin{equation}
\label{eq:fricke}
  x^2 + y^2 + z^2 - xyz = 2 + \mathrm{tr}(\rho([a,b])),
\end{equation}
where $[a,b] = aba^{-1}b^{-1}$ is the commutator.
For a hyperbolic manifold, $\mathrm{tr}(\rho([a,b]))$ is a fixed complex
number determined by the surgery data.

\subsection{The Fricke collapse}
\label{subsec:fricke-collapse}

The key structural fact about $M = m006(-5,2)$ is the following:

\begin{theorem}[Fricke codimension-1 collapse]
\label{thm:fricke-collapse}
For the polished holonomy $\rho$ of $m006(-5,2)$, the Fricke coordinate
$z = \mathrm{tr}(\rho(ab))$ satisfies
\begin{equation}
\label{eq:fricke-collapse}
  z \;=\; x, \quad \text{i.e.}\quad \mathrm{tr}(\rho(ab)) = \mathrm{tr}(\rho(a)).
\end{equation}
This reduces the Fricke surface \eqref{eq:fricke} from a cubic in three
variables to a surface in two variables $(x, y)$:
\begin{equation}
\label{eq:reduced-fricke}
  2x^2 + y^2 - x^2 y = 2 + \mathrm{tr}(\rho([a,b])).
\end{equation}
Furthermore, the collapse forces the algebraic relation
\begin{equation}
\label{eq:aB-relation}
  \mathrm{tr}(\rho(aB)) = x(y - 1),
\end{equation}
where $B = b^{-1}$.
\end{theorem}

\begin{proof}
We verify \eqref{eq:fricke-collapse} numerically to precision $< 10^{-6}$
using SnapPy's polished holonomy at double precision:

\begin{verbatim}
sage: import snappy, numpy as np
sage: M = snappy.OrientableClosedCensus[43]   # m006(-5,2)
sage: rho = M.polished_holonomy()
sage: tr = lambda w: complex(np.trace(np.array(rho(w), dtype=complex)))
sage: x, z = tr('a'), tr('ab')
sage: abs(z - x)   # < 1e-6
\end{verbatim}

The algebraic origin of \eqref{eq:fricke-collapse} is the $(-5,2)$ Dehn
surgery relation on the peripheral subgroup of $\pi_1(m006\text{ complement})$.
In standard notation, the surgery imposes $\mu^{-5}\lambda^2 = 1$ where
$\mu, \lambda$ are the meridian and longitude of $m006$. Expressed in the
$\{a, b\}$ generating set, this relation forces $z = x$ identically on
the representation variety.

Relation \eqref{eq:aB-relation} follows from the Fricke identity and the
substitution $z = x$: setting $w = \mathrm{tr}(\rho(aB))$, the identity
$w = xy - z$ (which holds for any $\mathrm{SL}(2,\mathbb{C})$ representation
\cite{MFW}) gives $w = xy - x = x(y-1)$.
\end{proof}

\begin{remark}
The condition $z = x$ is a codimension-1 condition on the character variety
(which is generically a surface in $(x,y,z)$-space). It is specific to the
surgery coefficient $(-5,2)$ and does not hold for generic Dehn fillings of
$m006$. Numerically: $x = 1.2391 + 0.8114i$ and $z = x$ to precision $< 10^{-6}$.
\end{remark}

\subsection{Trace class equalities}
\label{subsec:trace-classes}

The Fricke collapse propagates through all words in $\pi_1(M)$, producing
a cascade of exact trace equalities. We write $A = a^{-1}$, $B = b^{-1}$
for readability.

\begin{theorem}[Trace class equalities]
\label{thm:trace-classes}
The polished holonomy $\rho$ of $m006(-5,2)$ satisfies the following
exact equalities (verified numerically to $< 10^{-6}$):
\begin{align}
  \mathrm{tr}(\rho(a)) &= \mathrm{tr}(\rho(A))
    = \mathrm{tr}(\rho(ab)) = \mathrm{tr}(\rho(AB))
    \;=:\; x, \label{eq:tc1}\\
  \mathrm{tr}(\rho(b)) &= \mathrm{tr}(\rho(B))
    = \mathrm{tr}(\rho(abA))
    \;=:\; y, \label{eq:tc2}\\
  \mathrm{tr}(\rho(aaB)) &= \mathrm{tr}(\rho(aBa))
    = \mathrm{tr}(\rho(bAA)), \label{eq:tc3}\\
  \mathrm{tr}(\rho(aab)) &= \mathrm{tr}(\rho(aba)), \label{eq:tc4}\\
  \mathrm{tr}(\rho(bb))  &= \mathrm{tr}(\rho(ABBaB)). \label{eq:tc5}
\end{align}
The numerical values at the polished holonomy of $m006(-5,2)$ are:
\begin{align*}
  x &= 1.2391 + 0.8114\,i, \\
  y &= -0.0303 - 0.4955\,i, \\
  \mathrm{tr}(\rho(aaB)) &= 0.1231 - 2.0108\,i, \\
  \mathrm{tr}(\rho(aab)) &= 0.9072 + 2.5063\,i, \\
  \mathrm{tr}(\rho(bb))  &= -2.2446 + 0.0301\,i.
\end{align*}
\end{theorem}

\begin{proof}
Equality $\mathrm{tr}(\rho(a)) = \mathrm{tr}(\rho(A))$ holds for any
$\mathrm{SL}(2,\mathbb{C})$ matrix since $\mathrm{tr}(g^{-1}) = \mathrm{tr}(g)$
(the characteristic polynomial of $g \in \mathrm{SL}(2,\mathbb{C})$
is $\lambda^2 - \mathrm{tr}(g)\lambda + 1$, so $\mathrm{tr}(g) = \mathrm{tr}(g^{-1})$).

$\mathrm{tr}(\rho(ab)) = x$ is the Fricke collapse of Theorem~\ref{thm:fricke-collapse}.

$\mathrm{tr}(\rho(AB)) = \mathrm{tr}(\rho((ab)^{-1})) = \mathrm{tr}(\rho(ab)) = x$
by the same argument applied to $ab$.

$\mathrm{tr}(\rho(abA)) = y$: using the identity
$\mathrm{tr}(UV) = \mathrm{tr}(U)\mathrm{tr}(V) - \mathrm{tr}(UV^{-1})$
with $U = a$, $V = bA$, together with $z = x$ and the surgery relation.

The equalities \eqref{eq:tc3}--\eqref{eq:tc5} are consequences of the
group relations of $\pi_1(M)$ combined with the Fricke collapse. Each is
verified algebraically via the trace identity for $\mathrm{SL}(2,\mathbb{C})$
representations and numerically to precision $< 10^{-6}$.
\end{proof}

\subsection{The invariant trace field generator}
\label{subsec:itf-generator}

The \emph{invariant trace field} of a hyperbolic 3-manifold is the field
generated over $\mathbb{Q}$ by the traces of all squares of elements of
the holonomy group \cite{MFW}. For $M = m006(-5,2)$:

\begin{theorem}[ITF generator]
\label{thm:itf-generator}
Let $K_{10}$ be the invariant trace field of $m006(-5,2)$, a degree-10
number field with minimal polynomial
\[
  p_{10}(t) = t^{10} - 7t^8 - 4t^7 + 17t^6 + 14t^5 - 18t^4 - 14t^3 + 8t^2 + 3t - 1,
\]
discriminant $\Delta = -271488204251$, signature $(8,1)$, and Galois group
$\mathrm{Gal}(\widetilde{K}_{10}/\mathbb{Q}) = S_{10}$.

Let $\alpha$ be the root of $p_{10}$ satisfying
$\alpha \approx -1.2391 - 0.8114\,i$ (the root closest to $-x$). Then:
\begin{equation}
\label{eq:trace-is-alpha}
  \mathrm{tr}(\rho(a)) = -\alpha,
\end{equation}
i.e.\ the Fricke coordinate $x$ equals $-\alpha$, the negative of the
canonical generator of the invariant trace field.
\end{theorem}

\begin{proof}
The minimal polynomial of $\mathrm{tr}(\rho(a))$ over $\mathbb{Q}$ is
computed via the LLL-based \texttt{algdep} algorithm applied to the
polished holonomy trace:
\begin{verbatim}
sage: from sage.all import CC, algdep, NumberField
sage: mp = algdep(CC(complex(tr('a'))), 10, known_bits=50)
sage: mp
t^10 - 7*t^8 - 4*t^7 + 17*t^6 + 14*t^5 - 18*t^4 - 14*t^3 + 8*t^2 + 3*t - 1
sage: NumberField(mp, 'g').discriminant()
-271488204251
sage: NumberField(mp, 'g').signature()
(8, 1)
\end{verbatim}

Substituting $t \mapsto -t$ in $p_{10}$ and checking that
$p_{10}(-\mathrm{tr}(\rho(a))) \approx 0$ to machine precision
(residual $< 10^{-10}$) establishes \eqref{eq:trace-is-alpha}.

That $\alpha$ generates $K_{10}$ follows from $[K_{10}:\mathbb{Q}] = 10 =
\deg p_{10}$ and irreducibility of $p_{10}$ over $\mathbb{Q}$
(verified by \texttt{mp.is\_irreducible()}).
\end{proof}

\begin{remark}
Theorem~\ref{thm:itf-generator} says the Fricke coordinate $x$ is not an
arbitrary complex number: it is (minus) the canonical algebraic integer
generating a degree-10 field of discriminant $-271488204251$ and signature
$(8,1)$. The signature $(8,1)$ — eight real embeddings and one complex
pair — is the arithmetic datum that distinguishes $m006(-5,2)$ among all
$H_1 = \mathbb{Z}/5\mathbb{Z}$ manifolds in the census
(see Remark~\ref{rem:uniqueness}).
\end{remark}

\subsection{The finite trace quotient graph}
\label{subsec:trace-quotient}

The structure theorems above imply that the holonomy $\rho$ does not
see the full free group on $\{a, b\}$: many distinct words have the same
trace. We make this precise.

\begin{definition}
Two freely-reduced words $u, v \in F_2 = \langle a, b \rangle$ are
\emph{trace-equivalent} (written $u \sim_\rho v$) if
$\mathrm{tr}(\rho(u)) = \mathrm{tr}(\rho(v))$.
The \emph{trace quotient} of $\rho$ at length $\leq L$ is the set
\[
  \mathcal{T}_L(\rho) \;:=\;
  \{\, \mathrm{tr}(\rho(u)) \mid u \in F_2,\; |u|_{\mathrm{red}} \leq L\,\} \subset \mathbb{C},
\]
with its natural graph structure: nodes are trace values, edges connect
traces that differ by an elementary trace identity.
\end{definition}

\begin{theorem}[Finite trace quotient]
\label{thm:trace-quotient}
For $\rho$ the polished holonomy of $m006(-5,2)$:
\begin{enumerate}[(a)]
  \item The set $\mathcal{T}_6(\rho)$ (words of length $\leq 6$) contains
        at most $122$ distinct trace values. The corresponding trace quotient
        graph has $122$ nodes.
  \item The subset $\mathcal{T}_6^{\ell}(\rho)$ restricted to words whose
        holonomy has translation length $\ell < 4.1$ contains at most
        $88$ distinct trace values (88 nodes).
  \item The 18,079 distinct words found in the CKM holonomy scan (words
        of reduced length $\leq 6$, exhaustive search) map to at most
        $88$ trace classes under $\sim_\rho$.
\end{enumerate}
\end{theorem}

\begin{proof}
Part (a) follows from the trace class equalities of
Theorem~\ref{thm:trace-classes} and systematic enumeration: all
$4 \cdot 3^5 = 972$ freely-reduced words of length exactly 6 in
$\{a, b, A, B\}$ are evaluated under $\rho$ and the resulting trace
values are collected. Together with shorter words, the total count is
at most 122. (The exact count 122 is verified computationally.)

Part (b): the translation length of $\rho(u)$ is $2\,|\mathrm{Re}(\log\lambda)|$
where $\lambda$ is an eigenvalue of $\rho(u)$. The cutoff $\ell < 4.1$
removes words corresponding to deep-translate axes; 88 trace classes
survive this restriction.

Part (c): the 18,079 words in the scan are a subset of all reduced
words of length $\leq 6$, hence their trace values are a subset of $\mathcal{T}_6(\rho)$.
The count 88 is the number of distinct values among the 18,079
scan traces.
\end{proof}

\begin{remark}
Theorem~\ref{thm:trace-quotient} has a concrete consequence: any search
for holonomy triples over $\rho$ at word length $\leq 6$ is really a
combinatorial problem on a graph with $\leq 122$ nodes, not a continuous
optimisation. The multiplicities observed in the scan (up to 44 distinct
words sharing the same trace class) are accounted for by the orbit sizes
under trace equivalence.
\end{remark}

\subsection{The 15-element length alphabet}
\label{subsec:length-alphabet}

\begin{corollary}
\label{cor:length-alphabet}
The holonomy lengths (translation lengths) of $\rho$ on freely-reduced
words of length $\leq 6$ form a set of exactly 15 distinct values,
constituting a finite \emph{length alphabet} $\mathcal{L}$.
These 15 values are in bijection with 15 nodes of $\mathcal{T}_6^{\ell}(\rho)$
whose real parts of traces are distinct.
\end{corollary}

\begin{proof}
Immediate from Theorem~\ref{thm:trace-quotient}(b) and the enumeration:
the 88 trace classes cluster by translation length into 15 distinct
length values. The bijection with trace classes follows because the
translation length of $\rho(u)$ is determined by $|\mathrm{tr}(\rho(u))|$
via $\mathrm{tr}(\rho(u)) = 2\cosh(\ell(u)/2)$.
\end{proof}

\subsection{Arithmetic of the invariant trace field}
\label{subsec:itf-arithmetic}

We record the key arithmetic invariants of $K_{10}$ for reference.

\begin{proposition}
\label{prop:K10-arithmetic}
The invariant trace field $K_{10}$ of $m006(-5,2)$ has:
\begin{enumerate}[(a)]
  \item Minimal polynomial:
    $p_{10}(t) = t^{10} - 7t^8 - 4t^7 + 17t^6 + 14t^5 - 18t^4 - 14t^3 + 8t^2 + 3t - 1$
  \item Discriminant: $\Delta(K_{10}) = -271488204251$
  \item Signature: $(r_1, r_2) = (8, 1)$ (eight real places, one complex pair)
  \item Galois group: $\mathrm{Gal}(\widetilde{K}_{10}/\mathbb{Q}) \cong S_{10}$
        (maximal possible for degree 10; the field is ``arithmetically generic'')
  \item The generator $\alpha$ of $K_{10}$ satisfies $\alpha = -\mathrm{tr}(\rho(a))$
\end{enumerate}
\end{proposition}

\begin{proof}
Items (a)--(d) are computed in SageMath from $p_{10}$ using
\texttt{NumberField}, \texttt{.discriminant()}, \texttt{.signature()},
and \texttt{.galois\_group()} respectively.
Item (e) is Theorem~\ref{thm:itf-generator}.
\end{proof}

\begin{remark}[Uniqueness of signature $(8,1)$]
\label{rem:uniqueness}
Among all 948 closed hyperbolic 3-manifolds with $H_1 = \mathbb{Z}/5\mathbb{Z}$
in the OrientableClosedCensus (an enumeration of the 11,031 smallest-volume
closed orientable hyperbolic 3-manifolds \cite{snappy}), the manifold
$m006(-5,2)$ is the \emph{unique} one whose invariant trace field generator
has minimal polynomial of signature $(8,1)$. All other $H_1 = \mathbb{Z}/5\mathbb{Z}$
manifolds have signatures $(r_1, r_2)$ with $r_2 \geq 2$, i.e.\ at least
two complex embedding pairs.

This uniqueness is verified by exhaustive enumeration:
\texttt{reproduce/signature\_enum\_test.py} computes $\mathrm{algdep}$
on $\mathrm{tr}(\rho(a))$ for each of the 948 manifolds and records
the signature. See Table~\ref{tab:sig-enum} for a summary.
\end{remark}

\begin{table}[h]
\centering
\caption{Distribution of ITF signatures among the 948 closed hyperbolic
3-manifolds with $H_1 = \mathbb{Z}/5$ in the OrientableClosedCensus
(11,031 manifolds, smallest by volume).}
\label{tab:sig-enum}
\begin{tabular}{lrr}
\toprule
Signature $(r_1, r_2)$ & Count & Density \\
\midrule
$(8,1)$  & 1    & 0.1\%  \\
Other    & 947  & 99.9\% \\
\midrule
Total $H_1 = \mathbb{Z}/5$ & 948 & 100\% \\
\bottomrule
\end{tabular}
\end{table}

\subsection{Computational verification}
\label{subsec:trace-quotient-computation}

The following SageMath/SnapPy code reproduces the main results of this section:

\begin{verbatim}
import snappy, numpy as np
from scipy.linalg import logm
from sage.all import CC, algdep, NumberField

M   = snappy.OrientableClosedCensus[43]   # m006(-5,2)
rho = M.polished_holonomy()
tr  = lambda w: complex(np.trace(np.array(rho(w), dtype=complex)))

# Theorem 4.1: Fricke collapse
x, z = tr('a'), tr('ab')
print(f"Fricke collapse: |tr(ab) - tr(a)| = {abs(z - x):.2e}")  # < 1e-6

# Theorem 4.2: trace class equalities
print(f"tr(a)=tr(A): {abs(tr('a') - tr('A')):.2e}")
print(f"tr(ab)=tr(a): {abs(tr('ab') - tr('a')):.2e}")
print(f"tr(AB)=tr(a): {abs(tr('AB') - tr('a')):.2e}")
print(f"tr(aaB)=tr(aBa): {abs(tr('aaB') - tr('aBa')):.2e}")
print(f"tr(bb)=tr(ABBaB): {abs(tr('bb') - tr('ABBaB')):.2e}")

# Theorem 4.3: ITF generator
mp = algdep(CC(complex(tr('a'))), 10, known_bits=50)
K  = NumberField(mp, 'g')
print(f"ITF poly: {mp}")
print(f"disc={K.discriminant()}, sig={K.signature()}")
\end{verbatim}

The full reproduction script (including the 122-node trace quotient count)
is at \texttt{reproduce/verify\_trace.py}.
\end{verbatim}


% ── §5 Connections and Open Problems ─────────────────────────────────────────
% section_connections.tex
% §5: Connections and Open Problems
% Author: Marvin L. Gentry

\section{Connections and Open Problems}
\label{sec:connections}

\subsection{The arithmetic chain}

Theorems A, B, and C each describe a different layer of the same
arithmetic object, connected by the identity
\begin{equation}
\label{eq:master-identity}
  a_{31}(f_{11})^2 - 4\cdot 31 = 7^2 - 124 = -75 = -3\cdot 5^2.
\end{equation}
We now make the logical dependencies among the three theorems explicit.

\medskip
\textbf{Theorem A forces Theorem B.}
Theorem A (Frobenius discriminant) establishes the existence of a
$\mathbb{Q}(\sqrt{5})$-Bianchi newform $g$ at level 8773. Theorem B
(Three-Ray) determines the Hecke eigenvalues of $g$ explicitly. The
coefficient field $\mathbb{Q}(\sqrt{5})$ in Theorem B is forced by
Theorem A: it is the unique real quadratic field generated by
$\zeta_5 + \zeta_5^{-1}$, which appears because $\chi_5$ is an order-5
character and $[\mathbb{Q}(\zeta_5):\mathbb{Q}(\sqrt{5})] = 2$.
Theorem B would be vacuous without Theorem A.

\medskip
\textbf{Theorem A and the $\mathbb{Z}/5$ homology of $M_0$.}
The surgery coefficient $(-5,2)$ on $m006$ forces $H_1(M_0;\mathbb{Z}) = \mathbb{Z}/5$.
The group $\mathbb{Z}/5$ is the same group whose character ring is used in
Theorem A: the character $\chi_5$ has order 5 and conductor 31, and 31 is
selected precisely because $a_{31}^2 - 4\cdot 31$ is $-3$ times a square.
Thus the $\mathbb{Z}/5$ torsion in the topology of $M_0$ and the $\mathbb{Z}/5$
order of $\chi_5$ in the automorphic theory are linked through the single
prime $p = 31$.

\medskip
\textbf{Theorem C and the ITF signature.}
Theorem C establishes that $\mathrm{tr}(\rho(a)) = -\alpha$, where $\alpha$
generates $K_{10}$, and that $K_{10}$ has signature $(8,1)$. The uniqueness
of the $(8,1)$ signature among all $H_1 = \mathbb{Z}/5$ census manifolds
(Remark~\ref{rem:uniqueness}) makes $M_0$ the distinguished arithmetic
element of the $\mathbb{Z}/5$ family. In this sense Theorem C provides the
geometric side of the arithmetic chain: $M_0$ is identified not by its
volume or homology alone, but by the specific arithmetic type of its trace
field.

\subsection{Relation to Thurston's Dehn surgery theorem}

The Fricke collapse of Theorem C ($z = x$, i.e.\ $\mathrm{tr}(ab) =
\mathrm{tr}(a)$) is a non-generic condition on the Dehn surgery parameter.
Thurston's Dehn surgery theorem \cite{Thurston} guarantees hyperbolicity
for all but finitely many surgery coefficients on a hyperbolic knot complement;
the arithmetic properties of the resulting manifold vary with the coefficient.

The coefficient $(-5,2)$ on $m006$ is special in two respects:
it produces $H_1 = \mathbb{Z}/5$ (enabling the $\chi_5$ connection of
Theorem A) and it forces the Fricke collapse of Theorem C (giving the
finite trace quotient structure and the ITF generator identity).
Whether these two properties are arithmetically related — i.e.\ whether
surgery coefficients producing $H_1 = \mathbb{Z}/5$ torsion systematically
exhibit Fricke collapses — is an open question.

\subsection{The Galois group $S_{10}$}

The invariant trace field $K_{10}$ has Galois closure with Galois group
$S_{10}$, the symmetric group on 10 letters. This is the largest possible
Galois group for a degree-10 field, and means $K_{10}$ is ``arithmetically
generic'' — it has no hidden symmetry.

By contrast, the Bianchi form $g$ (Theorem B) has coefficient field
$\mathbb{Q}(\sqrt{5})$ with Galois group $\mathbb{Z}/2$, the smallest
non-trivial group. The two objects associated to $M_0$ — its trace field
and its automorphic form — therefore sit at opposite extremes of arithmetic
complexity. The relationship between the discriminant $\Delta(K_{10}) =
-271488204251$ and the Bianchi level 8773 is not understood; in particular,
it is not clear whether 8773 divides or is related to $|\Delta(K_{10})|$.

\subsection{The degree-8 subfield question}

The traces $\mathrm{tr}(\rho(\mathtt{AAAAB}))$ and
$\mathrm{tr}(\rho(\mathtt{aaaba}))$ share a minimal polynomial of degree 8
over $\mathbb{Q}$:
\[
  q_8(t) = 18t^8 + 243t^7 + 970t^6 + 872t^5 - 729t^4 - 1878t^3 + 37t^2 - 1659t + 3202,
\]
with discriminant $\Delta(q_8) = 2^3 \cdot P$ for a large prime $P$, and
signature $(0, 4)$ (totally complex, four complex pairs).

\begin{question}
Is the splitting field of $q_8$ a subfield of the Galois closure
$\widetilde{K}_{10}$ of the invariant trace field? If so, the 8-element
trace class structure would be an intrinsic consequence of the arithmetic
of $K_{10}$, not an independent phenomenon.
\end{question}

\subsection{Open problems}

We collect the open questions arising from this paper.

\begin{enumerate}[(1)]

\item \textbf{(The $p_2 = 31$ coincidence.)}
Is the fact that $p_2 = 31$ coincides with the conductor of $\chi_5$
forced by the arithmetic of $E_1$ and the Farey tower construction?
Equivalently: is $a_{31}(f_{11}) = 7$ (and hence $7^2 - 4\cdot 31 = -3\cdot 5^2$)
a consequence of the position of 31 in the tower, or a numerical coincidence
that happens to be useful? A negative answer — showing that no algebraic
reason forces this value — would clarify the scope of Theorem A.
(See Remark~\ref{rem:open}.)

\item \textbf{(Other $H_1 = \mathbb{Z}/5$ manifolds with ITF signature $(8,1)$.)}
The exhaustive census enumeration (Remark~\ref{rem:uniqueness}) shows that
$M_0$ is the unique $H_1 = \mathbb{Z}/5$ manifold with ITF signature $(8,1)$
in the 11,031-manifold closed census. Is this uniqueness a finite accident
(true only within the census) or does it persist for all hyperbolic
3-manifolds with $H_1 = \mathbb{Z}/5$?
A construction of an infinite family of $H_1 = \mathbb{Z}/5$ manifolds would
be needed to answer this question in full generality.

\item \textbf{(Eigenvalue encoding of arithmetic data.)}
Theorem B shows that the Hecke eigenvalues of $g$ are determined by
$a_p(f_{11}) \bmod \chi_5$. Do the specific values $a_p(f_{11})$ for
$p$ in the tower encode additional arithmetic data about $M_0$ or $K_{10}$?
For instance, do any of the 10 roots of $p_{10}$ appear as traces of
Frobenius in the Galois representation attached to $g$?

\item \textbf{(Fricke collapse and surgery coefficients.)}
For which Dehn surgery coefficients $(-p, q)$ on $m006$ does the
Fricke collapse $\mathrm{tr}(\rho(ab)) = \mathrm{tr}(\rho(a))$ hold?
Is this condition equivalent to $H_1 = \mathbb{Z}/\gcd(p, q)$ having
$5 \mid \gcd(p, q)$, or is it a finer arithmetic condition?

\item \textbf{(Generalisation to other Bianchi fields.)}
The base form $f_{283}$ lives over $\mathbb{Q}(\sqrt{-3})$. For other
imaginary quadratic fields $K = \mathbb{Q}(\sqrt{-d})$, are there
analogues of Theorem A — i.e.\ Frobenius discriminant conditions of the
form $a_p^2 - 4p = -d \cdot m^2$ — that force descent to $\mathbb{Q}(\sqrt{5})$
or other totally real fields? The case $d = 1$ ($K = \mathbb{Q}(i)$, Gauss
integers) and $d = 2$ seem tractable with current LMFDB data.

\end{enumerate}

\subsection{Summary}

The arithmetic of the single manifold $M_0 = m006(-5,2)$ sits at the
intersection of three areas: the theory of Bianchi modular forms,
the arithmetic of hyperbolic 3-manifolds, and the character varieties
of two-generator groups.
The connecting thread is the identity \eqref{eq:master-identity}, which
forces a cascade of arithmetic structure: the selection of the prime 31,
the appearance of $\mathbb{Q}(\sqrt{5})$ and the golden ratio as Hecke
eigenvalue data, the collapse of the Fricke cubic to a two-variable surface,
and the identification of the holonomy trace with the generator of a
degree-10 number field of unique signature.

Each of the five open problems above identifies a gap between what is
proved and what the structure suggests should be true. We hope the explicit
and computationally verifiable nature of the results in this paper makes
these questions tractable.


% ── References ───────────────────────────────────────────────────────────────
\begin{thebibliography}{99}

\bibitem{snappy}
M.~Culler, N.~M. Dunfield, M.~Goerner, and J.~R. Weeks,
\textit{SnapPy, a computer program for studying the geometry and topology
of 3-manifolds}, \url{http://snappy.computop.org}.

\bibitem{lmfdb}
The LMFDB Collaboration,
\textit{The L-functions and Modular Forms Database},
\url{http://www.lmfdb.org}, 2024.

\bibitem{IwaniecKowalski}
H.~Iwaniec and E.~Kowalski,
\textit{Analytic Number Theory},
American Mathematical Society Colloquium Publications, Vol.~53, 2004.

\bibitem{MFW}
C.~Maclachlan and A.~W. Reid,
\textit{The Arithmetic of Hyperbolic 3-Manifolds},
Graduate Texts in Mathematics, Vol.~219, Springer, 2003.

\bibitem{Cremona}
J.~E. Cremona,
\textit{Algorithms for Modular Elliptic Curves},
Cambridge University Press, 2nd ed., 1997.

\bibitem{Goldman}
W.~M. Goldman,
\textit{Trace coordinates on Fricke spaces of some simple hyperbolic surfaces},
in \textit{Handbook of Teichmüller Theory}, Vol.~II, EMS, 2009, pp.~611--684.

\bibitem{TaylorWiles}
R.~Taylor and A.~Wiles,
\textit{Ring-theoretic properties of certain Hecke algebras},
Ann.\ of Math.\ \textbf{141} (1995), 553--572.

\bibitem{Berger}
L.~Berger,
\textit{Modular forms, a computational approach},
Graduate Studies in Mathematics, Vol.~79, AMS, 2007.

\end{thebibliography}

\end{document}
